\section{Key Doctrinal and Pastoral Distinctions}

Each Catholic teaching below stands beside the twin distortions---by excess or confusion, and by defect or denial---that it excludes.

\begin{fourstable}{Topic}{Catholic Teaching}{Distortion by Excess / Confusion}{Distortion by Defect / Denial}
Last end and revelation & Natural reason sees that finite goods cannot satisfy the will; revelation discloses God's free ordination of humanity to the beatific vision (\emph{ST} I-II, q.\,2, a.\,8; \emph{Dei Filius}, ch.\ 2). & Treating natural desire as a right binding God, or claiming reason can reach beatitude unaided. & Treating finite goods, survival, or knowing God merely as First Cause as humanity's complete end.\\
New Law and commandment & Chiefly the grace of the Holy Spirit inscribed in the heart; secondarily the written Gospel, moral precepts, and sacraments (\emph{ST} I-II, qq.\,106--108). & Legalism reducing Christian life to external rule-keeping divorced from infused charity. & Antinomianism pitting the Spirit against Scripture, moral law, or visible sacraments.\\
Hope and its opposites & Infused hope leans on God's omnipotence and mercy to attain eternal beatitude (\emph{ST} II-II, q.\,17, aa.\,1--2). & Presumption: expecting pardon without repentance, glory without grace, or salvation by native power (\emph{ST} II-II, q.\,21, a.\,1). & Despair: judging falsely that God refuses pardon to the repentant or cannot heal one's sin (\emph{ST} II-II, q.\,20, a.\,1).\\
Petitionary prayer & Real secondary causality ordained eternally by immutable providence, aligning the heart while obtaining ordained effects (\emph{ST} II-II, q.\,83, a.\,2). & Treating prayer as a lever that alters God's will or guarantees chosen temporal outcomes. & Reducing prayer to self-talk or dismissing it as futile under divine omniscience.\\
Dryness and presence & Conserving presence, sanctifying indwelling, and sacramental presence remain distinct from sensible consolation (\emph{ST} I, q.\,8, a.\,3; q.\,43, a.\,3). & Measuring charity by warmth, or reading dryness as abandonment or lost grace. & Abandoning regular prayer and fidelity whenever sensible consolation fades.\\
Eucharist and reception & Sacrament of charity; worthy reception requires faith and freedom from conscious grave sin, with Reconciliation as ordinary remedy (\emph{ST} III, q.\,80, a.\,4). & Treating subjective desire as sufficient to receive Communion in conscious mortal sin. & Neglecting spiritual communion when deferred, or treating past venial sins as a bar.\\
Suffering and merit & Suffering is an evil; relieve removable pain; unite unavoidable suffering freely to Christ's complete Passion; merit belongs only to free acts in charity (\emph{ST} I-II, q.\,114). & Treating pain as good, meritorious, or personal punishment; self-harm; neglecting lawful remedies. & Denying that unavoidable suffering freely united to Christ in charity bears spiritual fruit.\\
Deification and union & Glory makes the blessed \emph{deiform} by created participation (2 Pet 1:4; \emph{ST} I, q.\,12, a.\,5), keeping personhood and divine incomprehensibility (\emph{ST} I, q.\,12, a.\,7). & Pantheistic absorption, identity of essence, erasure of self, or reversing ``God is charity'' (1 John 4:8, 16) into ``love is God.'' & Reducing eternal life to an external wage or distant admiration without immediate vision.\\
\end{fourstable}
